% Generated draft: requires mathematical and source review.
\documentclass[11pt]{article}
\usepackage[margin=25mm]{geometry}
\usepackage{fontspec}
\setmainfont{Latin Modern Roman}
\usepackage{amsmath,amssymb,mathtools,unicode-math}
\setmathfont{Latin Modern Math}
\usepackage{enumitem,needspace,xurl,hyperref}
\setlength{\parindent}{0pt}
\setlength{\parskip}{5pt}
\setlength{\emergencystretch}{3em}
\newcommand{\N}{\mathbb{N}}
\newcommand{\Z}{\mathbb{Z}}
\newcommand{\Q}{\mathbb{Q}}
\newcommand{\R}{\mathbb{R}}
\newcommand{\C}{\mathbb{C}}
\newcommand{\F}{\mathbb{F}}
\newcommand{\A}{\mathbb{A}}
\newcommand{\PP}{\mathbb{P}}
\newcommand{\E}{\mathbb{E}}
\newcommand{\eps}{\varepsilon}
\newcommand{\dd}{\,\mathrm{d}}
\begin{document}
\section*{UnsolvedMath ID: 30005778}
\textbf{Source problem code:} catalogue entry
\subsection{Definitions and mathematical statement}
Let \(n\geq 3\) and consider the group data specified in the supplied record:
\[
G=\prod_{i=1}^{n}\bigl(D_i\rtimes_{\varphi_i}M_i\bigr),
\qquad M=\prod_{i=1}^{n}M_i.
\]
Each factor \(D_i\rtimes_{\varphi_i}M_i\) is required to be a Frobenius group, each \(M_i\) is abelian of prime-power order, and \(|M|\) is a perfect square. No requirement that the primes associated with different \(M_i\) coincide is supplied. In the usual finite Frobenius-group interpretation, \(D_i\) is the Frobenius kernel and \(M_i\) is a complement whose nonidentity elements fix no nonidentity element of \(D_i\) under the specified action. The record does not explicitly state finiteness of the \(D_i\), although its surrounding discussion concerns finite-dimensional Hopf algebras.

For a field \(K\), the group algebra \(KG\) has coproduct \(\Delta(g)=g\otimes g\), counit \(\varepsilon(g)=1\), and antipode \(S(g)=g^{-1}\) for \(g\in G\). An invertible element \(J\in KG\otimes_K KG\) is a twist, with the convention in the supplied excerpt, if
\[
(1\otimes J)(\operatorname{id}\otimes\Delta)(J)
=(J\otimes1)(\Delta\otimes\operatorname{id})(J),
\qquad
(\varepsilon\otimes\operatorname{id})(J)
=(\operatorname{id}\otimes\varepsilon)(J)=1.
\]
The twisted Hopf algebra \((KG)^J\) has the same multiplication and counit as \(KG\), and coproduct
\[
\Delta^J(h)=J\Delta(h)J^{-1}.
\]
The record requires that \(J\) arise from a \(2\)-cocycle on the character group \(\widehat M\). For a specified character group and coefficient field, a normalized multiplicative \(2\)-cocycle is a map \(\alpha:\widehat M\times\widehat M\to K^\times\) satisfying
\[
\alpha(\chi,\psi)\alpha(\chi\psi,\rho)
=\alpha(\psi,\rho)\alpha(\chi,\psi\rho),
\qquad
\alpha(1,\chi)=\alpha(\chi,1)=1.
\]
However, the record does not specify the character field, the hypotheses on \(K\), the construction of \(J\), or any nondegeneracy condition on the cocycle. The square-order assumption is retained without treating it as a supplied nondegeneracy hypothesis.

Under the usual Hopf-algebra meaning, a Hopf algebra \(H\) is simple if its only normal Hopf subalgebras are \(K1\) and \(H\). Here a Hopf subalgebra \(A\subseteq H\) is normal when, for all \(h\in H\) and \(a\in A\), both
\[
\sum h_{(1)}aS(h_{(2)})\in A,
\qquad
\sum S(h_{(1)})ah_{(2)}\in A,
\quad
\Delta(h)=\sum h_{(1)}\otimes h_{(2)}.
\]
The supplied record does not explicitly define simplicity.

\textbf{Question as supplied, with unresolved qualifications.} For the above data and \(n\geq3\), is \((KG)^J\) simple? Read universally, this asks whether simplicity holds for every choice of the stated group data and every allowed cocycle twist. The missing field hypotheses, cocycle restrictions, finiteness convention, and intended meaning of simplicity must be resolved before this can be regarded as a fully specified conjecture. The supplied original excerpt introduces twisting and Hopf orders but does not itself contain this question.
\subsection{Short English statement}
The supplied question asks whether twisting the group algebra of a product of at least three specified Frobenius groups produces a simple Hopf algebra. The assumptions on the field and cocycle, and the intended quantification over twists, remain incomplete. The record labels the problem open; that label and its claimed source verification are supplied curation, not independent verification.
\subsection{Sources}
\begin{itemize}
\item Local UnsolvedMath registry, permanent ID \(30005778\), title \emph{Simplicity of Twisted Hopf Algebras}, source problem code \texttt{catalogue entry}. The supplied record labels the statement \texttt{corrected\_verified} and says it was reconstructed and checked against the official report. These are supplied metadata; no independent source check is asserted here.
\item \emph{Mini-Workshop: Bridging Number Theory and Nichols Algebras via Deformations} (2024), Oberwolfach Reports, pp.~235--272, DOI \texttt{10.4171/owr/2024/5}. Supplied source and literature link: \url{https://doi.org/10.4171/owr/2024/5}.
\item The supplied original excerpt identifies the contribution \emph{Non-existence of Hopf orders for certain Hopf algebras}, Elisabetta Masut, joint work with Giovanna Carnovale and Juan Cuadra, with a report page marker of 251. This attributes the excerpt; the record has no separately identified problem proposer.
\item Supplied additional literature source: \emph{Lifting via cocycle deformation}, \url{https://arxiv.org/abs/1212.5279}. Its precise relationship to the simplicity question is not explained in the record.
\item Literature assessment dated exactly \texttt{2026-08-22}: the supplied curation labels the problem open and states that finite-dimensional pointed Hopf algebras admitting Hopf orders satisfy a referenced conjectural conclusion, while the general statement is not a theorem and remaining cases are unresolved. This is supplied dated curation, not independently verified. The referenced conjecture is not stated, so this assessment cannot be identified with the simplicity question from the supplied evidence alone.
\item The original excerpt says that admitting a Hopf order over a number ring implies an unspecified conjecture and suggests Hopf orders may help classify semisimple Hopf algebras. It roughly describes a Hopf order as a Hopf algebra over a subring \(R\subseteq K\) whose scalar extension to \(K\) is the given Hopf algebra. The antecedent conjecture and the person referred to as having proved the implication are absent from the excerpt.
\end{itemize}
\end{document}
